% Cover-letter philosophy paragraph (~160w). Drop-in for BioSci's cover letter to Nature / Cell / Science editors.
% Emphasizes the conceptual contribution without over-claiming. Names REAL, LRS, RareShield.

Our work identifies and resolves a structural limit in single-cell batch-integration evaluation. The field's standard quality metrics --- ARI, NMI, ASW and their variants in scIB, scIB-E and RBET --- are measurable functionals of annotator labels. We prove that every such functional is invariant under integration deformations that can erase unannotated rare cell subpopulations (Theorem~1). The failure mode the community has worried about for at least five years therefore sits outside its own evaluation apparatus: absence of evidence has been systematically confused with evidence of absence. We replace the semantic criterion (``does the population have a name?'') with a syntactic one (``is there reproducible structure across independent observations?''), operationalised by a label-free detectability framework (REAL, with per-seed and per-method scores LRS($w$), LRS($f$); hierarchical-REAL extends the envelope by telescoping across compartment depth) calibrated on held-out known rare types under an explicit exchangeability bound, and paired with a protection regulariser (RareShield) that demonstrably reduces erasure at pan-cancer-atlas scale \emph{within the envelope defined by our minimax detectability bound} (detectable regime $\pi\Delta^{2}\geq 1.5\times 10^{-4}$ at $4.9\times 10^{6}$ cells under batch-heterogeneity penalty $\Lambda\approx 3$). The manuscript states the epistemic scope of the claim and the boundary conditions under which the bound fails, distinguishing detection from ontological discovery.
